\documentclass[12pt]{article}
\usepackage[utf8]{inputenc}
\usepackage{amsmath, amssymb, amsthm}
\usepackage{geometry}

% Set margins to match a standard homework document
\geometry{a4paper, margin=1in}

\title{Recitation 2}
\parskip=1ex
\author{Math 311 Fall 2026 \\ Rutgers University}
\date{}

\begin{document}

\maketitle

\noindent{} \textbf{Question 1:}
\begin{enumerate}
      \item [(a)] True. Consider $ B = {1} \subseteq R$.
            Then  \[
                  \inf(B) = 1 = \sup(B).
            \]
            Therefore,
            \[
                  \inf B \geq \sup B.
            \]
      \item [(b)] False, every non-empty finite set will contain both its supremum and infimum.
            Thus a finite set cannot contain its infimum but fail to contain its supremum.
      \item [(c)] Define
            \[
                  B = (0, 1] \cap \mathbb{Q}.
            \]
            Then $B$ is a bounded nonempty subset of $\mathbb{Q}$, where
            \[
                  \sup B = 1 \in B,
            \]
            and
            \[
                  \inf B = 0 \notin B.
            \]
\end{enumerate}

\noindent{} \textbf{Question 2:}
\begin{enumerate}
      \item [(a)] The axiom of completeness says that every nonempty set in $\mathbb{R}$ that is bounded above has a supremum.

            Similarly, for an ordered field $F$ to be complete, every non-empty subset $S$ of $F$ that is bounded above must have a supremum.
      \item [(b)]
            Let
            \[
                  S = \mathbb{Q} \cap (0, \pi).
            \]
            Suppose, for sake of contradiction, that $S$ has a supremum in the ambient ordered field of $\mathbb{Q}$, say
            \[ s = \sup S. \]
            Notice that $\pi > x$ for every $x \in S$, so
            \[
                  s \leq \pi
            \]
            Since $s \in \mathbb{Q}$ and $\pi \notin \mathbb{Q}$,
            \[
                  s < \pi
            \]
            Furthermore, $\mathbb{Q}$ is dense in $\mathbb{R}$, so there must exist some $q \in \mathbb{Q}$ such that
            \[
                  s < q < \pi.
            \]
            Since $s \in S$, it follows that $0 < s < \pi$. So
            \[
                  0 < q < \pi,
            \]
            and $q \in S$.
            But $q > s$, which contradicts our assumption that $s$ is an upper bound on $S$.
            Therefore, $S$ has no supremum in $\mathbb{Q}$.

            If we instead choose the ambient ordered field to be $\mathbb{R}$,
            then $\pi$ is still an upper bound on $S$ since every $x \in S$ satisfies
            \[
                  \pi \geq x.
            \]
            For sake of contradiction, suppose there exists a $u < \pi$ such that
            \[
                  u \geq x
            \]
            for every $x \in S$.
            Then, by density of $\mathbb{Q}$, there exists a $q \in \mathbb{Q}$ such that
            \[
                  u < q < \pi
            \]
            But, then $q \in S$, so $u$ cannot be an upper bound on $S$.
            Therefore $\pi$ must be the least upper bound.

            This proves that $\mathbb{Q}$ is not complete,
            because even though $S \subseteq \mathbb{Q}$ is nonempty and bounded above,
            it has no supremum in $\mathbb{Q}$
      \item [(c)]
            Let
            \[ Q_1 = \mathbb{Q} \cap (0, \pi),~ Q_2 = \mathbb{Q} \cap (\pi, 100)\]
            For every $a \in Q_1$
            \[
                  a < \pi
            \]
            and for every $b \in Q_2$
            \[
                  b > \pi.
            \]
            Thus,
            \[
                  a < \pi < b,
            \]
            so
            \[
                  a < b.
            \]

            Suppose, for contradiction, that there exists some $c \in \mathbb{Q}$ such that
            \[
                  a \leq c \leq b
            \]
            for every $a \in Q_1$ and every $b \in Q_2$.
            Since $c \in \mathbb{Q}$ and $\pi \notin \mathbb{Q}$, either $c < \pi$ or $c > \pi$.

            (Case 1) If $c < \pi$, then density of $\mathbb{Q}$ says there exists some $q \in \mathbb{Q}$ such that
            \[
                  c < q < \pi.
            \]
            So $q \in Q_1$.
            But $q > c$, contradicting $c \geq a$ for every $a \in Q_1$.

            (Case 2) If $c > \pi$, then density of $\mathbb{Q}$ says there exists some $q \in \mathbb{Q}$ such that
            \[
                  \pi < q < c.
            \]
            So $q \in Q_2$.
            But $q < c$, contradicting $c \leq b$ for every $b \in Q_2$.

            Since both cases are impossible, this $c$ cannot exist in $\mathbb{Q}$.
            Therefore,  there is no single $c \in \mathbb{Q}$ such that
            \[
                  a \leq c \leq b
            \]
            for every $a \in Q_1$ and every $b \in Q_2$.

\end{enumerate}

\noindent{} \textbf{Question 3:}
\begin{enumerate}
      \item [(a)] Let $X, Y \subseteq \mathbb{R}$ be nonempty. Assume that $X$ is bounded above, $Y$ is bounded below,
            \[
                  x \leq y, \text{ for every } x \in X \text{ and for every } y \in Y.
            \]
            That is, every $y$ is an upper bound on $X$.
            Since $\sup X$ is the least upper bound, it follows that
            \[
                  \sup X \leq y, \text{ for every } y \in Y.
            \]
            That is, $\sup X$ is a lower bound on $Y$.
            Since $\inf Y$ is the greatest lower bound, it follows that
            \[
                  \sup X \leq \inf Y.
            \]
      \item [(b)]
            Continuing from part (a), assume
            \[
                  X \cup Y = \mathbb{R}.
            \]
            From part (a), we have
            \[
                  \sup X \leq \inf Y.
            \]
            For sake of contradiction, suppose that the strict inequality holds. So
            \[
                  \sup X < \inf Y.
            \]
            Then, the midpoint would be
            \[
                  m = \frac{\sup X + \inf Y}{2}.
            \]
            So
            \[
                  \sup X < m < \inf Y
            \]
            Since $X \cup Y = \mathbb{R}$, either $m \in X$ or $m \in Y$.

            (Case 1) If $m \in X$, then by definition of supremum,
            \[
                  \sup X \geq m.
            \]
            However, this contradicts that $\sup X < m$.

            (Case 2) If $m \in Y$, then by definition of infimum,
            \[
                  \inf Y \leq m
            \]
            However, this contradicts that $\inf Y > m$.

            Both cases lead to contradictions.
            Since $m \in \mathbb{R}$ and $X \cup Y = \mathbb{R}$, this is impossible.
            The strict inequality does not hold, so from $\sup X \leq \inf Y$ we get
            \[
                  \sup X = \inf Y.
            \]
\end{enumerate}

\noindent{} \textbf{Question 4}

Let $A \subseteq \mathbb{R}$ satisfy the relative Nested Interval Property, and
$E \subseteq A$ be nonempty and bounded above by some element of $A$.
Define
\[
      U = \{u \in A : u \text{ is an upper bound for } E\}
\]
and
\[
      L = A \setminus U.
\]
\begin{enumerate}
      \item [(a)]
            Assume, for the sake of contradiction, that $U$ has no least element.
            By definition, $E$ is bounded above by some element of $A$.
            So there exists some $u \in A$ that is an upper bound on $E$.
            Then $u \in U$, and so
            \[ U \neq \emptyset. \]

            Take some $e \in E$.
            For sake of contradiction, suppose $e \in U$.
            Then by definition of $U$,
            \[
                  e \leq u \text{ for every } u \in U.
            \]
            That is, $e$ is the least element of $U$.
            However, this contradicts our assumption that $U$ has no least element.
            So $e \notin U$, which means $e \in L$. Thus,
            \[
                  L \neq \emptyset.
            \]

            Take some $\ell \in L$ and $u \in U$.
            Since $\ell \notin U$, $\ell$ is not an upper bound on $E$,
            there exists some $e \in E$ such that
            \[
                  e > \ell.
            \]
            Since $u$ is an upper bound on $E$,
            \[
                  u \geq e.
            \]
            Together,
            \[
                  \ell < e \leq u.
            \]
            Thus,
            \[
                  \ell < u \text{ for every } \ell \in L \text{ and } u \in U.
            \]

            Suppose, for sake of contradiction, that there exists some maximum $\ell_0$ in $L$.

            Since $\ell_0 \in L$, it is not an upper bound for $E$.
            So there exists some $e \in E$ such that
            \[
                  e > \ell_0.
            \]
            Since $\ell_0$ is the largest element of $L$ and $e$ is greater, $e \notin L$.
            Furthermore, because $e \in E \subseteq A$, it follows that $e \in U$.
            With $e \in U$ and $e \in E$, it follows that
            \[
                  e \leq u \text{ for every } u \in U.
            \]
            That is, $e$ is the least element of $U$.
            However, this contradicts our assumption that $U$ has no least element.
            Therefore, it must be that $L$ has no largest element.
      \item [(b)]
            Since $L$ is nonempty and bounded above, and $U$ is nonempty and bounded below,
            completeness of $\mathbb{R}$ gives
            \[
                  \alpha = \sup L
                  \qquad \text{and} \qquad
                  \beta = \inf U.
            \]
            Because $\alpha = \sup L$ and $L$ has no largest element, we may
            recursively choose $(a_n) \subseteq L$ so that
            \[
                  a_1 \leq a_2 \leq \cdots
            \]
            and
            \[
                  \alpha-\frac{1}{n}<a_n\leq\alpha.
            \]
            Similarly, because $\beta=\inf U$ and $U$ has no least element, we may
            choose $(b_n)\subseteq U$ so that
            \[
                  b_1\geq b_2\geq\cdots
            \]
            and
            \[
                  \beta\leq b_n<\beta+\frac{1}{n}.
            \]
      \item [(c)]
            Define
            \[
                  I_n = [a_n, b_n] \cap A.
            \]
            Recall from part (a) that
            \[
                  \ell < u \text{ for every } \ell \in L \text{ and every } u \in U.
            \]
            Since $(a_n) \subseteq L$ and $(b_n) \subseteq U$, it follows that
            \[
                  a_n < b_n \text{ for every } n \in \mathbb{N}.
            \]
            Furthermore, because $(a_n)$ is increasing while $(b_n)$ is decreasing, each interval will contain the next.
            That is,
            \[
                  I_1 \supseteq I_2 \supseteq I_3 \supseteq \cdots.
            \]

            Thus, by the relative Nested Interval Property,
            \[
                  \bigcap_{n=1}^{\infty} I_n \neq \emptyset.
            \]
            Choose some
            \[
                  x \in \bigcap_{n=1}^{\infty} I_n.
            \]
            Then for every $n \in \mathbb{N}$,
            \[
                  a_n \leq x \leq b_n.
            \]
            Suppose, for sake of contradiction, that $x < \alpha$. Then
            \[
                  \alpha - x > 0.
            \]
            By the Archimedean Property, there exists $n \in \mathbb{N}$ such that
            \[
                  \frac{1}{n} < \alpha - x.
            \]
            Thus,
            \[
                  x < \alpha - \frac{1}{n} < a_n.
            \]
            But, this contradicts $a_n \leq x$.
            Therefore,
            \[
                  x \geq \alpha.
            \]

            Similarly, suppose, for sake of contradiction, $x > \beta$
            Then
            \[
                  x - \beta > 0.
            \]
            By the Archimedean Property, there exists $n \in \mathbb{N}$ such that
            \[
                  \frac{1}{n} < x - \beta.
            \]
            Thus,
            \[
                  b_n < \frac{1}{n} + \beta < x.
            \]
            But, this contradicts $x \leq b_n$.
            Therefore,
            \[
                  x \leq \beta
            \]
            Together
            \[
                  \alpha \leq x \leq \beta.
            \]
            Hence,
            \[
                  x \in A \cap [\alpha,\beta].
            \]

            Since $A = L \cup U$, either $x \in L$ or $x \in U$.

            (Case 1) Say $x \in L$. Then since $\alpha$ is the supremum of $L$,
            \[
                  x \leq \alpha.
            \]
            But since $x \in [\alpha, \beta]$,
            \[
                  x \geq \alpha.
            \]
            So $x = \alpha$. Since $x\in L$, this means $\alpha\in L$.
            But this contradicts the fact that $L$ has no largest element.

            (Case 2) Say $x \in U$. Then since $\beta = \inf U$,
            \[
                  x \geq \beta.
            \]
            Since $x \in [\alpha, \beta]$,
            \[
                  x \leq \beta.
            \]
            So $x = \beta$. Since $x\in U$, this means $\beta\in U$, contradicting the fact that $U$ has no least element.

            Thus, neither case is possible.
            But $x\in A$, so one of them
            must hold. This is a contradiction.
            Hence, our original assumption that $U$ has no least element must be false.
            Therefore, $U$ has a least element, say, $s$.
            Since $U$ is the set of upper bounds of $E$,
            $s$ must be the least upper bound of $E$ in $A$.
            That is,
            \[
                  s = \sup_A E.
            \]

            Since $E \subseteq A$ was arbitrary, every nonempty subset of $A$
            that is bounded above by an element of $A$ has a least upper bound belonging to $A$.

            Therefore, by definition of completeness, we conclude that $A$ is complete.
\end{enumerate}

\noindent{} \textbf{Question 5}
\begin{enumerate}
      \item [(a)]
            Define $f: \mathbb{N} \cup \{0\} \rightarrow \mathbb{N}$ as $f(x) = x + 1$

            Injective: Suppose $f(y) = f(x)$.
            Then
            \[
                  x + 1 = y + 1.
            \]
            So
            \[
                  x = y.
            \]
            Thus, $f(x)$ is injective.

            Surjective:
            Let $y \in \mathbb{N}$.
            Choose $x = y - 1$. Then
            \[
                  f(x) = y - 1 + 1 = y.
            \]
            Thus, $f(x)$ is also surjective.

            Therefore $f(x)$ is bijective.

      \item [(b)]
            Define
            \[
                  f:\{0\}\cup D\to D
            \]
            as
            \[
                  f(x) =
                  \begin{cases}
                        \frac{1}{2}, \text{ if } x = 0 \\
                        \frac{1}{2^{n+1}}, \text{ if }x = \frac{1}{2^n}
                  \end{cases}
            \]

            Injectivity:

            (Case 1) Suppose $f(x)=f(y)$.

            If $f(x)=\frac12$, then the only possible input is $x=0$,
            and similarly $y=0$.

            (Case 2) If $x = \frac{1}{2^n}$.
            Suppose $f(x) = f(y)$.
            Then by definition of $f$,
            \[
                  x = y = \frac{1}{2^n}.
            \]

            Thus, $f$ is injective in all cases.

            Surjectivity:

            Let $y \in D$. Then $y = \frac{1}{2^n}$ for some $n \in \mathbb{N}$.
            If $n = 1$, then $y = \frac{1}{2}$. So we can just choose $x = 0$.
            Then
            \[
                  f(0) = \frac{1}{2} = y.
            \]

            If $n \geq 2$, then $y = \frac{1}{2^n}$. So we choose
            \[
                  x = \frac{1}{2^{n-1}}.
            \]
            Then
            \[
                  f(x) = \frac{1}{2^n} = y.
            \]

            Thus, $f$ is surjective.
            Therefore, $f$ is a bijective function.
      \item [(c)]
            Define
            \[
                  f:[0,1)\to(0,1)
            \]
            as
            \[
                  f(x)=
                  \begin{cases}
                        \frac12,         & x=0,                                         \\
                        \frac1{2^{n+1}}, & x=\frac1{2^n}\text{ for some }n\in\mathbb N, \\
                        x,               & \text{otherwise}.
                  \end{cases}
            \]
            Injectivity:

            (Case 1) Suppose $f(x)=f(y)$.

            If $f(x)=\frac12$, then the only possible input is $x=0$,
            and similarly $y=0$.

            (Case 2) If $x = \frac{1}{2^n}$.
            Suppose $f(x) = f(y)$.
            Then by definition of $f$,
            \[
                  x = y = \frac{1}{2^n}.
            \]
            (Case 3) For all the other $x$ values, suppose $f(x) = f(y)$.
            Then $x = y$.

            Thus, $f$ is injective in all cases.

            Surjectivity:

            Let $y \in (0,1)$.

            (Case 1) If $y = \frac{1}{2}$. Choose $x = 0$. Then
            \[
                  f(0) = y
            \]
            (Case 2) If $y = \frac{1}{2^n}$. Choose $x = \frac{1}{2^{n-1}}$.
            Then
            \[
                  f(x) = \frac{1}{2^{n -1 + 1}} = y.
            \]
            (Case 3) If $y$ is any other value in $(0,1)$, just choose $x = y$.
            Then
            \[
                  f(x) = y.
            \]
            Thus, $f$ is surjective.
            Therefore, $f$ is a bijective function.

      \item [(d)]
            Define
            \[
                  g:[0,1] \to [0,1)
            \]
            as
            \[
                  g(x)=
                  \begin{cases}
                        \frac{1}{2},           & x = 1,                                                 \\
                        1 - \frac{1}{2^{n+1}}, & x = 1 - \frac{1}{2^n} \text{ for some } n\in\mathbb N, \\
                        x,                     & \text{otherwise}.
                  \end{cases}
            \]
            Injectivity:
            (Case 1) Suppose $f(x)=f(y)$.

            If $f(x)=\frac12$, then the only possible input is $x=1$,
            and similarly $y=1$.

            (Case 2) If $y = 1 - \frac{1}{2^n}$. Suppose $g(x) = g(y)$.
            Then, by definition of $g$,
            \[
                  x = y = 1 - \frac{1}{2^n}
            \]
            (Case 3) For all the other $x$ values, suppose $g(x) = g(y)$.
            Then $x = y$.

            Thus, $g$ is injective in all cases.

            Surjectivity:

            Let $y \in [0,1)$.

            (Case 1) If $y = \frac{1}{2}$. Choose $x = 1$.
            Then
            \[
                  g(1) = \frac{1}{2} = y
            \]
            (Case 2) If $y = 1 - \frac{1}{2^n}$. Choose $x = 1 - \frac{1}{2^{n-1}}$.
            Then
            \[
                  g(x) = 1 - \frac{1}{2^{n -1 + 1}} = y.
            \]
            (Case 3) If $y$ is any other value in $[0,1)$, just choose $x = y$.
            Then
            \[
                  g(x) = y.
            \]
            Thus, $g$ is surjective.

            That is, $g$ is bijective.

            Recall from part (c), we concluded that the function
            \[
                  f: [0, 1) \to (0,1)
            \]
            is bijective as well.

            Hence, the composition of the two functions
            \[
                  f \circ g: [0, 1] \to (0, 1)
            \]
            must also be a bijection.
            Therefore,
            \[
                  [0, 1] \sim (0, 1)
            \]
\end{enumerate}


\end{document}

